\section{Complete cost checks}
\label{sec:cost-validation}
Correct work counts are necessary, but they do not make a latency prediction exact. Two loops
can visit the same number of rows while reading very different memory locations. A small
candidate set reused repeatedly can stay in cache; different candidates may require fresh
memory reads. The following checks distinguish these effects and evaluate complete C++ calls.

\subsection{Buffer sizes}
The standalone profile uses the shared C++ score table and top-$K$ code. Its path operation
allocates a root bitmap, makes both children at each split, keeps the alternatives alive,
releases each parent, enumerates the final leaf and releases the remaining masks. The check confirms that the surviving path and alternatives together contain every original
document exactly once.

Widths range from 512 bytes to 125 MB. The largest width is one bit over a billion document
positions. Thirteen stored planes and the live masks reached about 3.50 GB of process memory.
This measurement covers the specified buffers rather than the full 256-plane index.
Whole-index memory includes every stored plane, packed document and ID.

The larger buffers showed median split costs of roughly 0.34--0.37 ns per physical word in this
standalone loop. At 256 dimensions, contiguous scoring took about 27 ns per document at one
and four million rows. Small posting lists cost more per row because fewer rows share the work of preparation and
result selection. Compilation and the rest of the query still affect timing; these measured
costs do not apply to every CPU instruction or workload.

\subsection{Row access cost}
A second operation reproduces the leaf-scoring loop: read a bitmap word, find its set bits,
fetch those document codes and IDs, score them, and maintain top $K$. It uses the same scorer
as the sequential loop. One run reuses a candidate bitmap after warmup. Another prepares a
different bitmap before each measured trial.

\begin{center}\small
\begin{tabular}{@{}lrr@{}}\toprule
16M-document buffer & Reused candidates & Changing candidates\\\midrule
About 1,958 scored rows & 0.159 ms & 0.455 ms\\
About 62,500 scored rows & 4.400 ms & 8.271 ms\\\bottomrule
\end{tabular}\normalsize\end{center}
Both runs select the same fraction of rows, but select different IDs. Generating candidates
happens outside the timer, although it still changes which data the cache already holds. The difference shows why the scan's cost per row should
not be used directly for B's scattered leaf scores. It does not directly measure cache-miss
counters or establish fixed CPU frequency.

A three-term fit used a fixed cost, physical bitmap words and scored rows. It was fitted on
1M and 4M pools and checked on a 16M pool. The reused-candidate model underestimated one check
by 45.2\%. With changing candidates, the worst error was still an underestimate of 27.6\%.
Both fits remain saved. The largest row pool was about 640 MB; it was not a 40 GB row-store test.

The timer for fetching and scoring leaf rows already includes finding their set bits.
Adding the separate timer for finding those bits would count that work twice. Also, a fitted
fixed cost of zero does not mean score-table preparation is free: the fit may include its
cost in other terms.

\subsection{Larger C++ runs}
The native check uses 1M, 2M and 4M binary documents, then a separate 8M run. It uses the same
constructed distribution of strong and weak query coordinates. Packed codes are generated in small chunks and
the first million rows are checked against the earlier document pool.

An independent reference avoids expanding these rows into a large float matrix. For each query,
it counts strong and weak mismatches directly from the packed bits. With magnitudes $a$ and
$\epsilon$, the score is
\[
S=a(m-2e_{\rm strong})+\epsilon(d-m-2e_{\rm weak}).
\]
The reference sorts documents by this score and then ID. Small tests compare that calculation with the earlier direct
float sum, including a dimension count that does not fill the final packed word. All complete
C++ calls in the scaling check matched the independent reference.

For B, the leaf threshold grows with the corpus: $L=160N/10^6$. This preserves the intended
13 strong-coordinate cuts. With fixed depth and density, split words, leaf words and scored
rows all grow approximately in proportion to $N$. A simple complete-call model is then
\[
\widehat T(N)=a_0+b_0N.
\]
The fixed term covers work that stays the same as the corpus grows; the slope combines
work that grows with it. This model assumes the specified path and growing leaf threshold.
It does not justify keeping $L=128$ while increasing $N$ indefinitely.

\subsection{Model corrections}
The first fit minimized percentage error over individual queries but was used to predict a
median. Fast queries received more weight, which did not match the intended median prediction. For example, workloads
with times $[1,10,10]$ and $[2,20,20]$ have medians 10 and 20. Weighting each query by inverse
time favors the fast minority and predicts the wrong median. The corrected fit uses one median
per corpus size. A test fails for the mixed workload under the old calculation and stays correct
for a uniform workload. The original fit is preserved.

A linear row term was also insufficient for C's small postings. The shared result heap contributes
measurable selection cost when only a few hundred documents are scored. Under independently
ordered continuous scores, document $i$ enters the current top $K$ with probability $K/i$.
The expected number of replacements after the first $K$ rows is
\[
\sum_{i=K+1}^{F}\frac{K}{i}\approx K\ln(F/K).
\]
For fixed $K$, a useful approximation is
\begin{equation}
\widehat T_C(F)=a_C+b_CF+c_C\ln\!\bigl(\max(1,F/K)\bigr).
\label{eq:small-posting-time}
\end{equation}
Binary ties and key ordering change the exact replacement count, so the logarithm is a model
term, not an exact counter. A large contiguous run supplies the per-row cost; the native 1M/2M
measurements fit the remaining terms. This corrected the small-posting prediction without
changing the reference ranking or search implementation.

The revised 4M check helps diagnose the model, but its results had already been seen.
The fitted costs were fixed before the 8M run, which checks predictions on a larger corpus
with the same prescribed queries. It does not check performance on new queries.

\begin{center}\small
\begin{tabular}{@{}llrrr@{}}\toprule
Strong coordinates & Method & Measured ms & Predicted ms & Error\\\midrule
Fixed & A & 197.029 & 203.016 & $+3.0\%$\\
Fixed & B & 0.854 & 0.884 & $+3.6\%$\\
Fixed & C & 0.0481 & 0.0518 & $+7.7\%$\\
Changing & A & 196.664 & 197.816 & $+0.6\%$\\
Changing & B & 0.900 & 0.859 & $-4.5\%$\\
Changing & C & 202.767 & 207.086 & $+2.1\%$\\\bottomrule
\end{tabular}\normalsize\end{center}
All six checks met the earlier 20\% criterion, supporting the stated model over the checked
sizes. The largest observed error does not provide a guaranteed error bound at a billion
rows. The earlier failed fits remain part of the evidence.

\subsection{Leaf-size follow-up}
At one million rows, the ideal 13-bit group has mean $\lambda\approx122$ and standard deviation
about 11. A leaf size of 128 can force another cut when the strong-matching group is slightly
larger than average. A rough upper margin $\lambda+3\sqrt{\lambda}$ is about 155; choosing 160
avoids many unnecessary weak-coordinate cuts. The preceding 12-bit group has mean about 244,
so it will usually still be split. These are distribution-based settings to check, not universal
thresholds or a guarantee from a normal approximation.

The follow-up adds 16,384 clusters, neighboring key widths and leaf sizes, including
160 for the changing-coordinate case. It remeasures prior choices rather than comparing fresh
timings against old ones. The first 40 generated queries become the tuning set; query IDs
40--59 are used only after the new choices are frozen. Before a router is reused, the document and reference prefixes are checked against the earlier data.

At 1M rows and the 99\% target, the changing-coordinate follow-up selected B with one cluster
and $L=160$. It measured 0.1446 ms at 100\% binary recall. A measured 24.9235 ms at 99.95\%:
a paired speedup estimate of 172.3, with 95\% interval 169.1--177.2. C chose a 14-bit key and
measured 14.6301 ms at 100\%, but its paired speedup interval against A included one. That
point estimate alone does not establish a C advantage.

The fixed-coordinate follow-up selected a 16,384-cluster, two-probe scan and a global 13-bit
key index. Both reached 100\% recall: A measured 0.0140 ms and C 0.0128 ms. Their small timing gap again compares closely related ways to select candidates; it does
not prove that one method is always faster.

\subsection{Key width and median time}
When strong coordinates change with each query, a fixed $h$-coordinate key contains a random
number $S_h$ of the $m$ strong coordinates. The hypergeometric distribution describes this
count when coordinates are selected without repetition:
\begin{equation}
\Pr(S_h=s)=\frac{\binom ms\binom{d-m}{h-s}}{\binom dh}.
\label{eq:captured-strong}
\end{equation}
For $d=256$, $m=13$, the probability of capturing none is 52.75\% at $h=12$, 49.94\% at $h=13$,
and 47.27\% at $h=14$. A small width change can move a sample across the halfway point between
queries that score nearly all rows and queries that can discard about half.

Assuming the stated gap between strong and weak scores, and also that bounds from weak key
coordinates cannot rule out individual weak patterns, the work for a query capturing $s$ strong coordinates is approximately
\[
H=2^{h-s},\qquad F=N2^{-s}.
\]
The expected scored fraction is $\sum_s\Pr(S_h=s)2^{-s}$: about 71.30\% at $h=13$ and 69.44\%
at $h=14$. The mean work changes smoothly even when a small sample's median changes sharply.
The analysis therefore requires retaining uncertainty across queries and the full range of query times.
The condition is about the stated constructed distribution, not arbitrary embeddings or IVF.

\subsection{One-bit control}
The original real-data parameter grid started with four-bit keys. A final control gives C a one-bit key and
A's routing layout, with no candidate or lookup limit. It keeps A and B's frozen settings and
uses 100 further query IDs, excluding all 500 previously used IDs. No parameter is changed
based on this control's outcomes.

A measured 9.541 ms, B 9.855 ms and C 10.545 ms. All three reached 98.95\% binary recall,
narrowly below the requested 99\%. The report retains this miss without adding a tolerance
to classify it as a pass. There is no supported alternative speedup in this check. The original 200-query
evaluation remains a separate result, where its conservative choices met their sample targets.
A finite-sample selection rule does not guarantee the same mean recall on every later query sample.

\subsection{Text-query timing}
The final check runs actual text through the cached local Nomic model using Apple's MPS
backend, then applies
the same normalization, routes the resulting vector and calls the C++ search. Batch size is
one. Each sequence shape is warmed first; model loading, index building and reference creation
are outside the request timer. References are computed afterward from the actual encoded vectors.

\begin{center}\small
\begin{tabular}{@{}lrrrr@{}}\toprule
Method & Encode median & Retrieve median & Total median & Total p95\\\midrule
A & 6.488 ms & 9.406 ms & 16.313 ms & 19.201 ms\\
B & 6.500 ms & 9.651 ms & 16.303 ms & 19.217 ms\\
C & 7.543 ms & 11.201 ms & 18.912 ms & 21.898 ms\\\bottomrule
\end{tabular}\normalsize\end{center}
These are 20 existing independent-query IDs, three blocks and two repeats, with the earlier
frozen settings. All reached 99.45\% binary recall on this sample. The small A/B difference does
not establish a speedup. Each total is measured per request: its median is not the sum of the
two stage medians. The process lifetime peak was about 3.67 GB including reference construction.
The MPS driver reported about 1.08 GB allocated. Some of that may already be included in
the process memory figure (RSS), so adding the two values may double-count memory.

This text check applies to the real Nomic queries. The constructed strong-coordinate vectors
do not come from this encoder, so adding its time to those examples would not constitute a
measured end-to-end sparse-embedding pipeline.
